\section{Many contact scales and the contradiction}
\label{sec:contradiction}

We now apply Proposition~\ref{prop:contact-bound} to the opposed
arrangement.  An interval of depths containing no contact forces the first
contact above each positive cut in that interval to occur close to the
top.  The first-contact estimate makes this unlikely on average over the
cuts.  We first count scales of \emph{depth}; the finite mean common-block
width will then convert this count into the block-index count used in
Proposition~\ref{prop:entropy-upper}.

\begin{proposition}\label{prop:contact-lower}
Suppose that both coordinate signs have positive directional-transience
probability.  In the opposed arrangement of
Section~\ref{sec:arrangement}, let $m(J)$ be the expected number of
dyadic common-block index intervals $(2^j,2^{j+1}]$, $1\le j\le J$,
containing a departure-site contact, as in~\eqref{eq:contact-count}.
There are constants $c>0$ and $J_0<\infty$ such that
\[
 m(J)\ge \frac{cJ}{\log\log(3+J)}\qquad(J\ge J_0).
\]
\end{proposition}

\begin{proof}
All constants below may depend on the two slab laws.  Choose a large
integer $J$ and put
\begin{equation}\label{eq:lower-parameters}
 N=2^{2J},\qquad
 G=\left\lceil4\log_2\log(2+J)\right\rceil+3.
\end{equation}
For sufficiently large $J$, we have $3\le G<J$.  Depth means distance
below the initial top of the arrangement.
The choice $2^{-G}=O((\log J)^{-4})$ compensates for the logarithmic
losses below, while $G=O(\log\log J)$ leaves enough contact scales.

\paragraph{A bridge ending at a prescribed depth.}
Let $\mathcal C_N$ be the event that the positive tape has a cut at
depth $N$, and write
\[
 \mathbf Q_N=\mathbf P(\,\cdot\mid\mathcal C_N).
\]
The renewal identity gives
$\mathbf P(\mathcal C_N)=u_N^+\ge p_+>0$.
Under $\mathbf Q_N$, reverse the list of positive slabs through this cut
and translate its bottom to the origin.  Call the resulting chronological
path $Y$, and let $T_a$ denote its first hit of height $a$.
Lemma~\ref{lem:exact-bridge} says that $Y$ has law
$P_0(\,\cdot\mid T_N<T_{-1})$, stopped at $T_N$.
The steps within each slab retain their original order.  Denote the
lateral endpoint by $E$, so that the endpoint is $(N,E)$.  The negative
trajectory starts at $(N-1,E)$ and, relative to that starting point,
has the independent law $\widehat P_{-e_1}$.

The endpoint has a useful elementary tail bound.  On $\mathcal C_N$
there are at most $N$ positive slabs before depth $N$, since their
widths are positive integers.  The norm of their total lateral
displacement is at most the sum of their radii, hence at most the sum
of the first $N$ positive tape radii.  Proposition~\ref{prop:radius}
and Markov's inequality give
\begin{equation}\label{eq:lower-endpoint-tail}
 \EE_{\mathbf Q_N}|E|
 \le \frac{N E_+R}{p_+},
 \qquad
 \mathbf Q_N\{|E|>N^2\}\le \frac{C}{N}.
\end{equation}
This uses a spatial radius moment only.

\paragraph{An uncovered scale forces a late first contact.}
Call $l\ge0$ a contact depth label if a contact occurs at some integer
depth in $[2^l,2^{l+1})$.  For $G<l\le J$, let $U_l$ be the event
that none of the labels $l-G,l-G+1,\ldots,l$ is a contact depth label.
Equivalently, there is no contact at depths in
$[2^{l-G},2^{l+1})$.

Fix a positive cut depth
\[
 s\in I_l:=\bigl[2^l,\tfrac32\,2^l\bigr]\cap\ZZ,
 \qquad k=N-s.
\]
After the list reversal, this cut occurs at the bridge's first hit
$T_k$, and the remaining bridge stays at or above height $k$.
Indeed, every earlier slab has departure heights below $k$, while
every later slab has departure heights at least $k$.  This is a
property of the words, without a stopping-time assertion about the
cut.

There is a contact after $T_k$ and before $T_N$: the last positive
departure is $(N-1,E)$, the starting departure site of the negative
trajectory.  Let $\tau$ be the first contact since $T_k$, defined
pathwise against the entire negative trajectory.  Its depth is at
most $s<2^{l+1}$.  On $U_l$, it must therefore be strictly less than
$2^{l-G}$.  With
\[
 r=N-2^{l-G},
\]
the nearest-neighbor path must have reached height $r$ before this
contact.  Thus
\begin{equation}\label{eq:uncovered-first-contact}
 T_r\le\tau<T_N,
 \qquad Y_n\cdot e_1\ge k\quad(T_k\le n\le\tau).
\end{equation}
No restriction on drawdowns from the running maximum is needed.

For each \emph{deterministic} $s\in I_l$, let $\mathcal A_{l,s}$ be
the event in~\eqref{eq:uncovered-first-contact}, together with
$|E|\le N^2$, whether or not $s$ is a cut.  Proposition~
\ref{prop:contact-bound}, with upper stopping height $N$, applies
to this event.  Its width-probability difference satisfies
\begin{align}
 \Delta_{l,s}
 &=u^+_{s-2^{l-G}}-u^+_s\notag\\
 &\le 2^{l-G}F_+(s-2^{l-G})
 \le 2^{-G}a_l,
 \qquad a_l:=2^l F_+(2^{l-1}),
 \label{eq:lower-delta}
\end{align}
because $s-2^{l-G}\ge2^{l-1}$ and
Lemma~\ref{lem:renewal-tail} applies.  The same Lemma gives
\begin{equation}\label{eq:lower-summable-tail}
 A:=\sum_{l\ge1}a_l<\infty.
\end{equation}
Writing $\phi(x)=x\log^2(e/x)$ for $0<x\le1$ and $\phi(0)=0$,
the first-contact estimate is
\begin{equation}\label{eq:lower-fixed-cut-bound}
 \mathbf Q_N(\mathcal A_{l,s})
 \le C\log(2N)\,\phi(\Delta_{l,s}).
\end{equation}

\paragraph{Averaging over the cuts.}
Let $\mu_+=E_+L$.  The renewal-count strong law yields, almost surely
under the positive tape law,
\[
 \#\{\hbox{positive cuts in }I_l\}
 =\frac{2^{l-1}}{\mu_+}+o(2^l)
 \qquad(l\longrightarrow\infty).
\]
Fix $c_0=1/(4\mu_+)$, and let $\mathcal B_J$ be the event that
each interval $I_l$, $G<l\le J$, contains at least $c_0 2^l$
positive cuts.  Since $G=G(J)\to\infty$, the eventual almost-sure
estimate implies $\mathbf P(\mathcal B_J^c)=o(1)$.  Moreover,
\begin{equation}\label{eq:cut-density-conditioning}
 \mathbf Q_N(\mathcal B_J^c)
 \le \frac{\mathbf P(\mathcal B_J^c)}{p_+}=o(1).
\end{equation}
In particular, no independence between the cut at $N$ and the density
event is required.

On $\mathcal B_J\cap\{|E|\le N^2\}$, every uncovered label $l$
produces $\mathcal A_{l,s}$ for each actual cut $s\in I_l$.
Consequently,
\begin{equation}\label{eq:cut-averaging}
 \1_{U_l}\1_{\mathcal B_J}\1_{\{|E|\le N^2\}}
 \le \frac1{c_0 2^l}\sum_{s\in I_l}\1_{\mathcal A_{l,s}}.
\end{equation}
The right side has dropped the cut requirement on $s$.  This permits
the deterministic-barrier estimate~\eqref{eq:lower-fixed-cut-bound};
the number of summands is $O(2^l)$ and cancels the averaging factor.

We next check that the logarithms in this estimate do not require a
quantitative tail rate.  For $J\ge3$, put
\[
 B_J=\sum_{l=1}^J a_l\log_+^2(1/a_l),
 \qquad \log_+x:=\max\{\log x,0\},
\]
where a zero summand is interpreted as zero.  Terms with
$a_l\ge J^{-2}$ contribute at most $4A\log^2 J$.  On
$(0,J^{-2})$, the function $x\log^2(1/x)$ is increasing, so the
remaining at most $J$ terms contribute at most $4J^{-1}\log^2 J$.
Therefore
\begin{equation}\label{eq:weighted-tail-log}
 B_J\le 4(A+J^{-1})\log^2 J.
\end{equation}

For sufficiently large $J$, $2^{-G}a_l\le e^{-1}$ uniformly in
$l$, since $a_l\le A$.  The function $\phi$ is increasing on
$[0,e^{-1}]$.  Thus~\eqref{eq:lower-delta} gives, uniformly over
$s\in I_l$,
\[
 \phi(\Delta_{l,s})\le C2^{-G}a_l
 \bigl(G+1+\log_+(1/a_l)\bigr)^2.
\]
Taking expectations in~\eqref{eq:cut-averaging}, then using
\eqref{eq:lower-endpoint-tail},
\eqref{eq:lower-fixed-cut-bound}, and
\eqref{eq:cut-density-conditioning}, gives
\begin{align}
 \sum_{l=G+1}^J\mathbf Q_N(U_l)
 &\le J\mathbf Q_N(\mathcal B_J^c)+\frac{CJ}{N}
   +C\log(2N)\,2^{-G}\bigl[A(G+1)^2+B_J\bigr]\notag\\
 &=o(J).
 \label{eq:few-uncovered}
\end{align}
Indeed, $\log(2N)=O(J)$, while the choice~
\eqref{eq:lower-parameters} gives
$2^{-G}=O(\log^{-4}(2+J))$.  By~\eqref{eq:weighted-tail-log},
the last term is $O(J/\log^2(2+J))$.

\paragraph{Removing the conditioning.}
Let $H_J$ count contact depth labels in $\{1,\ldots,J\}$.
Equation~\eqref{eq:few-uncovered} and Markov's inequality imply
that with $\mathbf Q_N$-probability tending to one, at most $J/4$
of the tested labels are uncovered.  Since $G=o(J)$, at least
$J/2$ are then covered.  A single contact depth label $h$ can cover
only the labels $h,h+1,\ldots,h+G$.  Hence
\begin{equation}\label{eq:height-count-conditioned}
 \mathbf Q_N\left\{H_J\ge\frac{J}{2(G+1)}\right\}
 \longrightarrow1.
\end{equation}
For large $J$, all depths counted here are less than $2^{J+1}<N$.
On $\mathcal C_N$ they are therefore fully represented in the
positive list ending at depth $N$.  Undoing the translation and list
reversal leaves exactly the contacts in the original opposed
arrangement.  Thus~\eqref{eq:height-count-conditioned} implies
\begin{equation}\label{eq:height-count-unconditioned}
 \mathbf P\left\{H_J\ge\frac{J}{2(G+1)}\right\}
 \ge p_+-o(1).
\end{equation}
Only a probability bounded away from zero is asserted after removing
the conditioning.  This is enough for an expectation lower bound.

\paragraph{From depths to common-block indices.}
Write $K_1,K_2,\ldots$ for the common-block widths and
$W_i=K_1+\cdots+K_i$, with $W_0=0$.  By
Lemma~\ref{lem:common-blocks}, $\mu:=\mathbf E K_1<\infty$ and
$W_i/i\to\mu$ almost surely.  In particular, eventually
\[
 \frac\mu2 i\le W_i\le2\mu i.
\]
An integer depth $z\ge1$ belongs to the unique block
\[
 i(z)=\min\{i\ge1:W_i\ge z\},
 \qquad W_{i(z)-1}<z\le W_{i(z)}.
\]
The departure-height convention makes this true even at a common
cut depth.  For every sufficiently large $z$, the strong-law bounds
give
\begin{equation}\label{eq:depth-index-comparison}
 \frac{z}{2\mu}\le i(z)<\frac{2z}{\mu}+1.
\end{equation}
The finite initial blocks, however large their widths, disappear once
$z$ exceeds their cumulative width.

Define the dyadic index label of a block $i>1$ by
$j(i)=\lceil\log_2 i\rceil-1$, so that
$2^{j(i)}<i\le2^{j(i)+1}$.  It follows from
\eqref{eq:depth-index-comparison} that there is a deterministic
integer $c_1$, depending only on $\mu$, such that eventually
\begin{equation}\label{eq:bounded-label-shift}
 |j(i(z))-l|\le c_1
 \quad\hbox{whenever }2^l\le z<2^{l+1}.
\end{equation}
Almost surely this holds above a finite random depth.  Therefore
the probability that it holds for every contact with
$\lceil\sqrt J\rceil\le l\le J$ tends to one.

Intersect this event with the event in
\eqref{eq:height-count-unconditioned}.  The intersection still has
probability at least $p_+-o(1)$.  Discarding the first
$\lceil\sqrt J\rceil$ depth labels loses
$o(J/(G+1))$ labels.  Each remaining contact depth label supplies a
contact block whose index label lies within $c_1$ of it.  Conversely,
one index label can account for at most $2c_1+1$ depth labels.
Thus, on the intersection, at least $c_2J/(G+1)$ of the index labels
$1,\ldots,J+c_1$ contain a contact, for a fixed $c_2>0$ and large
$J$.  Taking expectations yields
\[
 m(J+c_1)\ge\frac{c_3J}{G+1}.
\]
Since $G+1=O(\log\log(3+J))$, replacing $J+c_1$ by the argument
of $m$ proves the proposition.
\end{proof}

\begin{proof}[Proof of Theorem~\ref{thm:main}]
Fix any nonzero real direction $\ell$.  Suppose that
$0<P_0(A_\ell)<1$.  Lemma~\ref{lem:sign-union} implies that
$P_0(A_{-\ell})>0$.  The spatial-radius estimate and
Proposition~\ref{prop:coordinate-reduction} then supply a coordinate
direction in which both signs have positive probability.  Reflecting
and relabeling that coordinate gives the setting of
Propositions~\ref{prop:entropy-upper} and~\ref{prop:contact-lower}.
They imply simultaneously
\[
 m(J)=O\!\left(\frac{J}{\log J}\right)
 \quad\hbox{and}\quad
 m(J)\ge\frac{cJ}{\log\log(3+J)}
\]
for all sufficiently large $J$, which is impossible.  Thus
$P_0(A_\ell)\in\{0,1\}$.

The reduction applies to every fixed real $\ell$ without rational
approximation.  Strict ellipticity supplied the fresh-row probabilities;
the spatial-radius moment was derived under coexistence.  Neither a
moment of a slab's duration nor a speed assumption entered the argument.
\end{proof}
